\section{Experimental Setup}
\label{sec:experiments}

\subsection{Experimental Questions}

Our experiments address three questions:

\begin{enumerate}
    \item \textbf{Framework validity}: Can both matrix-level and token-level objectives be implemented and trained in a unified framework? (H-E1)
    \item \textbf{Teacher context limits}: Does Phi-1.5 exhibit extrapolation artifacts at extended sequence lengths? (H-M1)
    \item \textbf{Representation stability}: Do token-level and matrix-level objectives produce different drift patterns across sequence lengths? (H-M2)
\end{enumerate}

\subsection{Models}

\textbf{Teacher}: Phi-1.5 (microsoft/phi-1\_5), 1.3B parameters. Trained on 2048-token sequences with learned absolute positional embeddings. We discovered that Phi-1.5 has a \emph{hard} 2048 token limit---attempting to process longer sequences produces IndexError on position embeddings, not degraded attention patterns. This architectural constraint defines the scope boundary for our length extrapolation analysis.

\textbf{Student}: Phi-Mamba with MOHAWK modifications:
\begin{itemize}
    \item Multi-head SSM structure matching attention heads
    \item Removed $\Delta$ parameter (open gates)
    \item Embedding and output layers initialized from teacher
\end{itemize}

\subsection{Datasets and Metrics}

\textbf{Training}: C4 dataset (allenai/c4) via streaming access. 1M tokens (PoC) / 1.5B tokens per condition (full). Phi tokenizer for consistent vocabulary.

\textbf{Drift analysis}: 500 samples per sequence length from C4 validation split. Lengths: 512, 1024, 1536, 2048 tokens.

\textbf{Primary metric}: Hidden state drift slope (linear regression of L2 drift against sequence length). Lower slope = more stable representations across lengths.

\textbf{Secondary metrics}: Drift ratio (Drift(2048) / Drift(512)), cosine similarity degradation, per-layer drift heatmaps.

\subsection{Hypotheses Tested}

\begin{table}[h]
\centering
\begin{tabular}{llll}
\toprule
\textbf{ID} & \textbf{Hypothesis} & \textbf{Gate} & \textbf{Success Criterion} \\
\midrule
H-E1 & Unified framework validates both objectives & MUST\_WORK & No execution errors \\
H-M1 & Phi-1.5 attention extrapolation artifacts & MUST\_WORK & Measurable degradation >2048 \\
H-M2 & Token-level drift slope < matrix-level & MUST\_WORK & CAB slope < MOHAWK slope (p<0.05) \\
H-M3 & F1 retention interaction effect & MUST\_WORK & Interaction p < 0.05 \\
\bottomrule
\end{tabular}
\caption{Hypotheses tested in this work.}
\label{tab:hypotheses}
\end{table}
